%************************************************
\chapter{Biochemical Formula Tracking \& Atomic Conservation}
\label{ch:biochemical-formulas}
%************************************************

\section{Introduction}
\label{sec:formula-tracking-intro}

Biochemical reactions are governed by fundamental conservation laws:
\begin{itemize}
    \item \textbf{Mass conservation}: Total mass remains constant
    \item \textbf{Elemental conservation}: Atoms of each element (\ce{C}, \ce{H}, \ce{O}, \ce{N}, \ce{P}, \ce{S}) are neither created nor destroyed
    \item \textbf{Charge conservation}: Net charge remains constant (in ionic reactions)
    \item \textbf{Energy conservation}: Total energy conserved (though often dissipated as heat)
\end{itemize}

Traditional modeling approaches often \textbf{neglect elemental conservation}:
\begin{itemize}
    \item \textbf{Abstract tokens}: Classical Petri nets use tokens without chemical identity
    \item \textbf{Implicit stoichiometry}: ODE models assume correct stoichiometry but don't verify
    \item \textbf{Error-prone}: Easy to introduce imbalanced reactions (missing cofactors, wrong coefficients)
\end{itemize}

\paragraph{This chapter presents} the fourth core innovation: \textbf{biochemical formula tracking} ($\rho$ component of Extended Bio-PN).

\subsection{Key Contributions}
\label{subsec:formula-key-contributions}

\begin{enumerate}
    \item \textbf{Elemental composition map}: $\rho: P \to \text{BiochemicalFormula}$ assigns formulas to places
    \item \textbf{Reaction formulas}: $\rho: T \to \text{ReactionFormula}$ specifies elemental transformations
    \item \textbf{Balance verification}: Automatic checking that $\sum(\text{atoms}_{\text{in}}) = \sum(\text{atoms}_{\text{out}})$ for each element
    \item \textbf{Elemental balance matrix}: $S_e$ captures elemental stoichiometry (rows = elements, cols = transitions)
    \item \textbf{Debugging aid}: Imbalanced reactions flag modeling errors
\end{enumerate}

\paragraph{Benefits:}
\begin{itemize}
    \item \textbf{Correctness}: Ensures biochemical validity (no phantom atoms)
    \item \textbf{Clarity}: Explicit chemical transformations visible in model
    \item \textbf{Automation}: Tools can auto-suggest cofactors if imbalance detected
    \item \textbf{Integration with databases}: Parse formulas from KEGG, ChEBI, PubChem
\end{itemize}

%************************************************
\section{Biochemical Formula Representation}
\label{sec:formula-representation}
%************************************************

\subsection{Hill Notation}
\label{subsec:hill-notation}

\textbf{Standard representation}: Chemical formulas use \textbf{Hill notation}:
\begin{itemize}
    \item Carbon (\ce{C}) first, then hydrogen (\ce{H}), then other elements alphabetically
    \item Example: \ce{C6H12O6} (glucose)
\end{itemize}

\paragraph{Format:} \texttt{C<count>H<count>N<count>O<count>P<count>S<count>}
\begin{itemize}
    \item Omit element if count = 0
    \item Omit count if = 1 (e.g., \ce{CH4}, not \ce{C1H4})
\end{itemize}

\paragraph{Examples:}
\begin{table}[htbp]
\centering
\caption{Common metabolite formulas in Hill notation}
\label{tab:hill-notation-examples}
\begin{tabular}{ll}
\toprule
\textbf{Compound} & \textbf{Formula} \\
\midrule
Glucose & \ce{C6H12O6} \\
ATP & \ce{C10H16N5O13P3} \\
Water & \ce{H2O} (special case: no carbon, \ce{H} first) \\
Ammonia & \ce{H3N} (no carbon, \ce{H} first, then alphabetical) \\
Pyruvate & \ce{C3H4O3} (deprotonated form at pH 7) \\
Lactate & \ce{C3H5O3} (deprotonated form) \\
\bottomrule
\end{tabular}
\end{table}

\subsection{Protonation States}
\label{subsec:protonation-states}

\paragraph{Challenge:} Metabolites exist in multiple protonation states depending on pH.

\paragraph{Example: Phosphate}
\begin{align*}
\text{pH} < 2 &: \ce{H3PO4} \quad \text{(phosphoric acid)} \\
\text{pH } 2\text{--}7 &: \ce{H2PO4-} \quad \text{(dihydrogen phosphate)} \\
\text{pH } 7\text{--}12 &: \ce{HPO4^{2-}} \quad \text{(hydrogen phosphate)} \\
\text{pH} > 12 &: \ce{PO4^{3-}} \quad \text{(phosphate)}
\end{align*}

\paragraph{Convention:} Use \textbf{physiological pH} (pH 7.0--7.4 for cytoplasm):
\begin{itemize}
    \item \textbf{ATP}: \ce{C10H12N5O13P3^{4-}} $\to$ Simplified as \ce{C10H16N5O13P3} (add 4\ce{H+} to neutralize)
    \item \textbf{ADP}: \ce{C10H12N5O10P2^{3-}} $\to$ Simplified as \ce{C10H15N5O10P2} (add 3\ce{H+})
    \item \textbf{Glucose-6-phosphate}: \ce{C6H11O9P^{2-}} $\to$ Simplified as \ce{C6H13O9P}
\end{itemize}

\paragraph{Approach in Extended Bio-PN:}
\begin{itemize}
    \item \textbf{Store formulas at physiological pH} (majority species)
    \item \textbf{Track \ce{H+} explicitly} if proton balance critical (e.g., oxidative phosphorylation, pH regulation)
    \item \textbf{Otherwise}: Assume protons equilibrate rapidly with bulk water (large \ce{H2O} pool buffers pH)
\end{itemize}

\subsection{Elemental Decomposition}
\label{subsec:elemental-decomposition}

\textbf{Internal representation}: Parse Hill notation into \textbf{element $\to$ count map}.

\paragraph{Data structure:}
\begin{equation}
\text{BiochemicalFormula} = \{\text{Element} \to \mathbb{N}\}
\end{equation}

where $\text{Element} \in \{\ce{C}, \ce{H}, \ce{O}, \ce{N}, \ce{P}, \ce{S}, \ldots\}$

\paragraph{Examples:}
\begin{align}
\text{Glucose} &: \{\ce{C} \to 6, \ce{H} \to 12, \ce{O} \to 6\} \\
\text{ATP} &: \{\ce{C} \to 10, \ce{H} \to 16, \ce{N} \to 5, \ce{O} \to 13, \ce{P} \to 3\}
\end{align}

\begin{algorithm}[htbp]
\caption{Parse Hill notation to elemental map}
\label{alg:parse-formula}
\begin{algorithmic}[1]
\Require Hill notation string $s$ (e.g., ``\ce{C6H12O6}'')
\Ensure Elemental map $\rho = \{\text{element} \to \text{count}\}$
\State $\rho \gets \{\}$, $i \gets 0$
\While{$i < |s|$}
    \State $\text{element} \gets s[i]$ \Comment{First character (uppercase)}
    \State $i \gets i + 1$
    \If{$i < |s|$ \textbf{and} $s[i]$ is lowercase}
        \State $\text{element} \gets \text{element} + s[i]$ \Comment{Two-letter element (e.g., \ce{Ca})}
        \State $i \gets i + 1$
    \EndIf
    \State $\text{count\_str} \gets $ ``''
    \While{$i < |s|$ \textbf{and} $s[i]$ is digit}
        \State $\text{count\_str} \gets \text{count\_str} + s[i]$
        \State $i \gets i + 1$
    \EndWhile
    \State $\text{count} \gets \text{int}(\text{count\_str})$ if $\text{count\_str} \neq $ ``'' else $1$
    \State $\rho[\text{element}] \gets \text{count}$
\EndWhile
\State \Return $\rho$
\end{algorithmic}
\end{algorithm}

\subsection{Formula Arithmetic}
\label{subsec:formula-arithmetic}

\subsubsection{Addition (Combining Molecules)}

For two formulas $\rho_1, \rho_2$, their sum is:
\begin{equation}
(\rho_1 + \rho_2)[e] = \rho_1[e] + \rho_2[e] \quad \forall e \in \text{Elements}
\end{equation}

\paragraph{Example:} Glucose + ATP
\begin{align}
\rho_{\text{Glc}} &= \{\ce{C}: 6, \ce{H}: 12, \ce{O}: 6\} \\
\rho_{\text{ATP}} &= \{\ce{C}: 10, \ce{H}: 16, \ce{N}: 5, \ce{O}: 13, \ce{P}: 3\} \\
\rho_{\text{Glc}} + \rho_{\text{ATP}} &= \{\ce{C}: 16, \ce{H}: 28, \ce{O}: 19, \ce{N}: 5, \ce{P}: 3\}
\end{align}

\subsubsection{Subtraction (Removing Molecules)}

For two formulas $\rho_1, \rho_2$, their difference is:
\begin{equation}
(\rho_1 - \rho_2)[e] = \rho_1[e] - \rho_2[e] \quad \forall e \in \text{Elements}
\end{equation}

Elements with count 0 are omitted from the result.

\subsubsection{Multiplication (Stoichiometric Scaling)}

For formula $\rho$ and coefficient $k \in \mathbb{N}$:
\begin{equation}
(k \cdot \rho)[e] = k \cdot \rho[e] \quad \forall e \in \text{Elements}
\end{equation}

\paragraph{Example:} 2 ATP
\begin{equation}
2 \cdot \{\ce{C}: 10, \ce{H}: 16, \ce{N}: 5, \ce{O}: 13, \ce{P}: 3\} = \{\ce{C}: 20, \ce{H}: 32, \ce{N}: 10, \ce{O}: 26, \ce{P}: 6\}
\end{equation}

%************************************************
\section{Reaction Formula Specification}
\label{sec:reaction-formulas}
%************************************************

\subsection{Format}
\label{subsec:reaction-format}

\paragraph{Reaction formula:} String specifying elemental transformation:
\begin{equation}
\langle\text{reactants}\rangle \to \langle\text{products}\rangle
\end{equation}

\paragraph{Reactants/Products:} Space-separated list of $\langle\text{stoichiometry}\rangle\langle\text{formula}\rangle$
\begin{itemize}
    \item Stoichiometry optional if = 1
    \item Separate multiple compounds with ``\texttt{+}''
\end{itemize}

\subsection{Examples}
\label{subsec:reaction-examples}

\subsubsection{Hexokinase (Glucose Phosphorylation)}

\begin{equation}
\ce{C6H12O6 + C10H16N5O13P3 -> C6H13O9P + C10H15N5O10P2 + H}
\end{equation}

\noindent Readable: \ce{Glucose + ATP -> G6P + ADP + H+}

\subsubsection{Phosphoglucose Isomerase (G6P $\rightleftharpoons$ F6P)}

\begin{equation}
\ce{C6H13O9P -> C6H13O9P}
\end{equation}

\noindent\textbf{Note}: G6P and F6P have the \textbf{same elemental composition} (\ce{C6H13O9P}) but different \textbf{structures} (aldose vs.\ ketose).
\begin{itemize}
    \item \textbf{Elemental balance}: \ce{C6H13O9P -> C6H13O9P} \checkmark\ (trivial)
    \item \textbf{Structural difference}: Not captured by elemental formulas (would need SMILES or InChI)
\end{itemize}

\subsubsection{Phosphofructokinase (F6P Phosphorylation)}

\begin{equation}
\ce{C6H13O9P + C10H16N5O13P3 -> C6H14O12P2 + C10H15N5O10P2 + H}
\end{equation}

\noindent Readable: \ce{F6P + ATP -> F-1,6-BP + ADP + H+}

\subsubsection{Lactate Dehydrogenase (Pyruvate Reduction)}

\begin{equation}
\ce{C3H4O3 + C21H29N7O14P2 -> C3H6O3 + C21H27N7O14P2}
\end{equation}

\noindent Readable: \ce{Pyruvate + NADH -> Lactate + NAD+}

\noindent\textbf{Note}: NADH (reduced) has 2 more \ce{H} than \ce{NAD+} (oxidized)

\subsection{Extended Syntax with Cofactors}
\label{subsec:cofactor-syntax}

\paragraph{Common pattern:} Many reactions require \textbf{cofactors} not explicitly shown in simplified notation.

\paragraph{Example: ATP Hydrolysis}
\begin{itemize}
    \item \textbf{Simplified}: \ce{ATP -> ADP + Pi}
    \item \textbf{Complete}: \ce{ATP + H2O -> ADP + Pi + H+}
\end{itemize}

\paragraph{Water addition:}
\begin{itemize}
    \item Hydrolysis reactions consume water
    \item Condensation reactions produce water
    \item Often omitted (implicit) because $[\ce{H2O}]$ is high and constant ($\sim$55~M in aqueous solution)
\end{itemize}

\paragraph{Proton release:}
\begin{itemize}
    \item Phosphorylation reactions often release \ce{H+}
    \item pH-dependent (buffered in vivo)
    \item Should track if modeling pH changes
\end{itemize}

\paragraph{Extended Bio-PN approach:}
\begin{enumerate}
    \item \textbf{Store complete reaction formulas} (including \ce{H2O}, \ce{H+}, cofactors)
    \item \textbf{Optionally omit high-concentration species} from network topology (\ce{H2O} place not drawn)
    \item \textbf{Verify balance including all species} (even if some places hidden)
\end{enumerate}

%************************************************
\section{Elemental Balance Verification}
\label{sec:elemental-balance}
%************************************************

\subsection{Balance Check Algorithm}
\label{subsec:balance-check}

For a single transition $t$ with reaction formula $\rho(t)$:

\begin{algorithm}[htbp]
\caption{Verify elemental balance for transition}
\label{alg:verify-balance}
\begin{algorithmic}[1]
\Require Transition $t$ with reaction formula $\rho(t)$
\Ensure Imbalance map $\Delta: \text{Element} \to \mathbb{Z}$ where $\Delta[e] = \text{output}_e - \text{input}_e$
\State Parse $\rho(t)$ into reactants and products
\State $\text{input\_atoms} \gets \{\}$, $\text{output\_atoms} \gets \{\}$
\For{each $(k, \rho_i)$ in reactants} \Comment{$k$ = stoichiometry, $\rho_i$ = formula}
    \For{each element $e \in \rho_i$}
        \State $\text{input\_atoms}[e] \gets \text{input\_atoms}[e] + k \cdot \rho_i[e]$
    \EndFor
\EndFor
\For{each $(k, \rho_j)$ in products}
    \For{each element $e \in \rho_j$}
        \State $\text{output\_atoms}[e] \gets \text{output\_atoms}[e] + k \cdot \rho_j[e]$
    \EndFor
\EndFor
\State $\text{all\_elements} \gets \text{input\_atoms}.\text{keys}() \cup \text{output\_atoms}.\text{keys}()$
\State $\Delta \gets \{\}$
\For{each element $e \in \text{all\_elements}$}
    \State $\text{input\_count} \gets \text{input\_atoms}[e]$ (default 0 if not present)
    \State $\text{output\_count} \gets \text{output\_atoms}[e]$ (default 0 if not present)
    \If{$\text{input\_count} \neq \text{output\_count}$}
        \State $\Delta[e] \gets \text{output\_count} - \text{input\_count}$
    \EndIf
\EndFor
\State \Return $\Delta$ \Comment{Empty if balanced}
\end{algorithmic}
\end{algorithm}

\subsection{Example: Hexokinase Balance Check}
\label{subsec:hexokinase-balance}

\paragraph{Reaction:} \ce{C6H12O6 + C10H16N5O13P3 -> C6H13O9P + C10H15N5O10P2 + H}

\paragraph{Input atoms:}
\begin{align}
\ce{C} &: 6 + 10 = 16 \\
\ce{H} &: 12 + 16 = 28 \\
\ce{O} &: 6 + 13 = 19 \\
\ce{N} &: 0 + 5 = 5 \\
\ce{P} &: 0 + 3 = 3
\end{align}

\paragraph{Output atoms:}
\begin{align}
\ce{C} &: 6 + 10 + 0 = 16 \\
\ce{H} &: 13 + 15 + 1 = 29 \\
\ce{O} &: 9 + 10 + 0 = 19 \\
\ce{N} &: 0 + 5 + 0 = 5 \\
\ce{P} &: 1 + 2 + 0 = 3
\end{align}

\paragraph{Imbalance:} $\Delta[\ce{H}] = +1$

\paragraph{Issue detected:} Hydrogen imbalance! This suggests using \textbf{ionic formulas} (actual protonation states) is more accurate.

\subsection{Handling Protonation Consistently}
\label{subsec:protonation-handling}

\subsubsection{Approach 1: Ionic Formulas (Charge-Accurate)}

\begin{itemize}
    \item Store formulas in actual protonation states (e.g., \ce{ATP^{4-}}: \ce{C10H12N5O13P3})
    \item Track charge as additional ``element''
    \item Balance both elements AND charge
    \item \textbf{Pro}: Biochemically accurate
    \item \textbf{Con}: Requires tracking pH, buffer capacity
\end{itemize}

\paragraph{Example: Hexokinase with ionic formulas}
\begin{equation}
\ce{C6H12O6 + C10H12N5O13P3^{4-} -> C6H11O9P^{2-} + C10H12N5O10P2^{3-} + H+}
\end{equation}

\paragraph{Balance check:}
\begin{table}[htbp]
\centering
\caption{Elemental balance with ionic formulas}
\label{tab:ionic-balance}
\begin{tabular}{lrrr}
\toprule
\textbf{Element} & \textbf{Input} & \textbf{Output} & \textbf{Balance} \\
\midrule
\ce{C} & $6 + 10 = 16$ & $6 + 10 + 0 = 16$ & \checkmark \\
\ce{H} & $12 + 12 = 24$ & $11 + 12 + 1 = 24$ & \checkmark \\
\ce{O} & $6 + 13 = 19$ & $9 + 10 + 0 = 19$ & \checkmark \\
\ce{N} & $0 + 5 = 5$ & $0 + 5 + 0 = 5$ & \checkmark \\
\ce{P} & $0 + 3 = 3$ & $1 + 2 + 0 = 3$ & \checkmark \\
Charge & $0 + (-4) = -4$ & $(-2) + (-3) + (+1) = -4$ & \checkmark \\
\bottomrule
\end{tabular}
\end{table}

\subsubsection{Approach 2: Charge-Neutral Formulas (Proton-Adjusted)}

\begin{itemize}
    \item Add/remove \ce{H+} to neutralize charges (e.g., ATP: \ce{C10H16N5O13P3})
    \item Ignore charge balance (assume bulk water buffers)
    \item \textbf{Pro}: Simpler (no charge tracking)
    \item \textbf{Con}: Hydrogen balance may appear off (protons exchanged with water)
\end{itemize}

\subsubsection{Recommended Approach: Hybrid}

\begin{enumerate}
    \item Use \textbf{ionic formulas} by default (biochemically accurate)
    \item \textbf{Track \ce{H+} explicitly} as a place (if pH regulation modeled)
    \item \textbf{OR}: Omit \ce{H+} tracking (assume buffered) and accept minor \ce{H} imbalances
\end{enumerate}

\paragraph{Extended Bio-PN implementation:}
\begin{itemize}
    \item Store formulas in \textbf{database-canonical form} (as retrieved from KEGG, ChEBI)
    \item Flag for each place: \texttt{track\_protonation} (boolean)
    \item If \texttt{track\_protonation=True}: Use ionic formulas, verify \ce{H} balance strictly
    \item If \texttt{track\_protonation=False}: Relax \ce{H} balance check (allow $\pm$ few \ce{H} due to proton exchange)
\end{itemize}

%************************************************
\section{Elemental Balance Matrix}
\label{sec:balance-matrix}
%************************************************

\subsection{Definition}
\label{subsec:balance-matrix-def}

\paragraph{Stoichiometric matrix} $S \in \mathbb{Z}^{|P| \times |T|}$: 
\begin{equation}
S[p,t] = \text{net token change of place } p \text{ by transition } t
\end{equation}

\paragraph{Elemental balance matrix} $S_e \in \mathbb{Z}^{|\text{Elements}| \times |T|}$:
\begin{equation}
S_e[e,t] = \text{net atom count of element } e \text{ by transition } t
\end{equation}

\subsection{Construction Algorithm}
\label{subsec:balance-matrix-construction}

\begin{algorithm}[htbp]
\caption{Construct elemental balance matrix}
\label{alg:construct-balance-matrix}
\begin{algorithmic}[1]
\Require Extended Bio-PN with places $P$, transitions $T$, formulas $\rho$
\Ensure Elemental balance matrix $S_e$
\State $\text{Elements} \gets \{\ce{C}, \ce{H}, \ce{O}, \ce{N}, \ce{P}, \ce{S}\}$
\State Initialize $S_e[e,t] \gets 0$ for all $e \in \text{Elements}$, $t \in T$
\For{each transition $t \in T$}
    \For{each element $e \in \text{Elements}$}
        \State $\text{input\_atoms}_e \gets \sum_{p \in {}^\bullet t} W(p,t) \cdot \rho(p)[e]$
        \State $\text{output\_atoms}_e \gets \sum_{p \in t^\bullet} W(t,p) \cdot \rho(p)[e]$
        \State $S_e[e, t] \gets \text{output\_atoms}_e - \text{input\_atoms}_e$
    \EndFor
\EndFor
\State \Return $S_e$
\end{algorithmic}
\end{algorithm}

\paragraph{Meaning:} $S_e[e,t]$ is the net production (+) or consumption (--) of element $e$ by transition $t$.

\subsection{Steady-State Condition}
\label{subsec:steady-state-condition}

At metabolic steady state, elemental balance requires:
\begin{equation}
S_e \cdot v = 0
\label{eq:steady-state-balance}
\end{equation}
where $v \in \mathbb{R}_{\geq 0}^{|T|}$ is the flux vector (firing rates of transitions).

\paragraph{Interpretation:} Sum of elemental production/consumption across all active reactions = 0 (conservation).

\subsection{Example: Simplified Glycolysis}
\label{subsec:glycolysis-balance-matrix}

\paragraph{Transitions:}
\begin{itemize}
    \item $t_1$: Hexokinase (\ce{Glucose + ATP -> G6P + ADP})
    \item $t_2$: PFK (\ce{F6P + ATP -> F-1,6-BP + ADP})
    \item $t_3$: Pyruvate kinase (\ce{PEP + ADP -> Pyruvate + ATP})
\end{itemize}

\paragraph{Elemental balance matrix} $S_e$ (6 elements $\times$ 3 transitions):

\begin{table}[htbp]
\centering
\caption{Elemental balance matrix for simplified glycolysis}
\label{tab:glycolysis-balance-matrix}
\begin{tabular}{lrrr}
\toprule
\textbf{Element} & $t_1$ (Hexokinase) & $t_2$ (PFK) & $t_3$ (PK) \\
\midrule
\ce{C} & 0 & 0 & 0 \\
\ce{H} & +1 & +1 & --1 \\
\ce{O} & 0 & +3 & --3 \\
\ce{N} & 0 & 0 & 0 \\
\ce{P} & 0 & +1 & --1 \\
\ce{S} & 0 & 0 & 0 \\
\bottomrule
\end{tabular}
\end{table}

\paragraph{Interpretation:}
\begin{itemize}
    \item \textbf{Carbon}: Conserved in all reactions ($S_e[\ce{C}, :] = 0$)
    \item \textbf{Hydrogen}: $t_1$ and $t_2$ release \ce{H+} (proton production), $t_3$ consumes \ce{H+}
    \item \textbf{Phosphorus}: $t_2$ adds phosphate, $t_3$ removes phosphate (net transfer from ATP to metabolites)
\end{itemize}

\paragraph{Steady-state flux:} Suppose $v = [v_1, v_2, v_3]^T$ (flux vector).

Elemental balance: $S_e \cdot v = 0$
\begin{align}
\text{For } \ce{C} &: 0 \cdot v_1 + 0 \cdot v_2 + 0 \cdot v_3 = 0 \quad \checkmark \text{ (always satisfied)} \\
\text{For } \ce{H} &: 1 \cdot v_1 + 1 \cdot v_2 - 1 \cdot v_3 = 0 \quad \Rightarrow \quad v_3 = v_1 + v_2 \\
\text{For } \ce{P} &: 0 \cdot v_1 + 1 \cdot v_2 - 1 \cdot v_3 = 0 \quad \Rightarrow \quad v_3 = v_2
\end{align}

\paragraph{Contradiction:} $v_3 = v_1 + v_2$ (from \ce{H} balance) vs.\ $v_3 = v_2$ (from \ce{P} balance).

\paragraph{Resolution:} This indicates the \textbf{model is incomplete}. Missing reactions:
\begin{itemize}
    \item Proton-producing steps between $t_2$ and $t_3$ (e.g., glyceraldehyde-3-phosphate dehydrogenase)
    \item Or: Protons exchanged with water (bulk buffer, not tracked)
\end{itemize}

\paragraph{Conclusion:} Elemental balance matrix reveals \textbf{missing reactions} or \textbf{implicit buffering}.

\subsection{Null Space Analysis}
\label{subsec:null-space-analysis}

\paragraph{Mathematical insight:} The \textbf{null space} of $S_e$ contains flux vectors that conserve all elements.

\begin{equation}
\mathcal{N}(S_e) = \{v \in \mathbb{R}^{|T|} : S_e \cdot v = 0\}
\end{equation}

\paragraph{Biological interpretation:} Feasible steady-state flux distributions.

\paragraph{Example (complete glycolysis, 10 reactions):}
\begin{itemize}
    \item $\text{rank}(S_e) = 5$ (5 elements: \ce{C}, \ce{H}, \ce{O}, \ce{N}, \ce{P} tracked; \ce{S} omitted if no sulfur-containing metabolites)
    \item $\text{nullity} = 10 - 5 = 5$ (5 degrees of freedom)
    \item \textbf{Meaning}: 5 independent flux modes conserve elemental balance
\end{itemize}

\paragraph{Applications:}
\begin{enumerate}
    \item \textbf{Flux balance analysis (FBA)}: Optimize flux $v$ subject to $S_e \cdot v = 0$ (elemental constraints)
    \item \textbf{Elementary flux modes (EFMs)}: Basis vectors for $\mathcal{N}(S_e)$ (minimal pathways)
    \item \textbf{Metabolic pathway inference}: Which flux distributions conserve mass?
\end{enumerate}

%************************************************
\section{Integration with Databases}
\label{sec:database-integration}
%************************************************

\subsection{KEGG Compound Formulas}
\label{subsec:kegg-integration}

\textbf{KEGG} (Kyoto Encyclopedia of Genes and Genomes) provides chemical formulas for $\sim$18,000 compounds.

\paragraph{Example API query:}
\begin{verbatim}
GET http://rest.kegg.jp/get/C00031
\end{verbatim}

\paragraph{Response (excerpt):}
\begin{verbatim}
ENTRY       C00031                      Compound
NAME        D-Glucose;
            Grape sugar;
            Dextrose
FORMULA     C6H12O6
EXACT_MASS  180.0634
MOL_WEIGHT  180.156
\end{verbatim}

\paragraph{Extended Bio-PN integration:}
\begin{itemize}
    \item Place annotation: \texttt{kegg\_id = "C00031"} (Glucose)
    \item Automatic formula lookup: $\rho(\text{Glucose}) = \texttt{fetch\_kegg\_formula("C00031")}$
    \item Cache formulas locally (avoid repeated queries)
\end{itemize}

\subsection{ChEBI Formulas}
\label{subsec:chebi-integration}

\textbf{ChEBI} (Chemical Entities of Biological Interest) provides more detailed chemical information.

\paragraph{Example:}
\begin{itemize}
    \item ChEBI:17234 (D-glucose)
    \item Formula: \ce{C6H12O6}
    \item Charge: 0
    \item Protonation state: Neutral
\end{itemize}

\paragraph{Advantage over KEGG:} ChEBI includes \textbf{charge information} and \textbf{protonation states}.

\paragraph{Example: ATP}
\begin{itemize}
    \item ChEBI:30616 (\ce{ATP^{4-}})
    \item Formula: \ce{C10H12N5O13P3}
    \item Charge: --4
\end{itemize}

\subsection{PubChem Formulas}
\label{subsec:pubchem-integration}

\textbf{PubChem} (NIH database, $>$100 million compounds) also provides formulas.

\paragraph{Example:}
\begin{itemize}
    \item PubChem CID 5793 (ATP)
    \item Molecular Formula: \ce{C10H16N5O13P3} (neutral form)
\end{itemize}

\paragraph{Recommendation:} Use \textbf{KEGG} for metabolites (curated, biologically relevant), \textbf{ChEBI} for detailed chemical properties, \textbf{PubChem} for broader chemical space.

%************************************************
\section{Automatic Cofactor Suggestion}
\label{sec:cofactor-suggestion}
%************************************************

\subsection{Detecting Imbalance}
\label{subsec:detect-imbalance}

\paragraph{Scenario:} User specifies reaction without all cofactors.

\paragraph{Example:}
\begin{equation}
\text{User input: } \ce{Glucose -> G6P}
\end{equation}

Parsed formula: \ce{C6H12O6 -> C6H13O9P}

\paragraph{Imbalance:}
\begin{align}
\ce{C} &: 6 \to 6 \quad \checkmark \\
\ce{H} &: 12 \to 13 \quad \text{(--1 \ce{H} missing in reactants, +1 \ce{H} in products)} \\
\ce{O} &: 6 \to 9 \quad \text{(--3 \ce{O} missing in reactants)} \\
\ce{P} &: 0 \to 1 \quad \text{(--1 \ce{P} missing in reactants)}
\end{align}

\paragraph{Interpretation:} Reaction requires adding atoms (\ce{H}, \ce{O}, \ce{P}).

\paragraph{Possible cofactors:}
\begin{itemize}
    \item \textbf{ATP}: Provides \ce{P} (+3), \ce{O} (+13), but wrong \ce{H}/\ce{O} ratio
    \item \textbf{Phosphate} (\ce{Pi}): Provides \ce{P} (+1), \ce{O} (+4), \ce{H} (+2)
    \item \textbf{Water}: Provides \ce{H} (+2), \ce{O} (+1)
\end{itemize}

\subsection{Cofactor Matching Algorithm}
\label{subsec:cofactor-matching}

\begin{algorithm}[htbp]
\caption{Suggest cofactors to balance imbalance}
\label{alg:suggest-cofactors}
\begin{algorithmic}[1]
\Require Imbalance map $\Delta: \text{Element} \to \mathbb{Z}$ where $\Delta[e] = \text{output} - \text{input}$
\Ensure List of suggested cofactor names (sorted by helpfulness)
\State $\text{cofactors} \gets \{$\ce{ATP}, \ce{ADP}, \ce{AMP}, \ce{Pi}, \ce{H2O}, \ce{NAD+}, \ce{NADH}, \ce{CoA}, \ce{CO2}$\}$
\State $\text{suggestions} \gets []$
\For{each cofactor $c \in \text{cofactors}$}
    \State $\Delta' \gets \Delta$ \Comment{Copy imbalance}
    \For{each element $e \in \rho(c)$}
        \State $\Delta'[e] \gets \Delta'[e] - \rho(c)[e]$
    \EndFor
    \State $\text{old\_score} \gets \sum_{e} |\Delta[e]|$
    \State $\text{new\_score} \gets \sum_{e} |\Delta'[e]|$
    \If{$\text{new\_score} < \text{old\_score}$}
        \State $\text{suggestions}.\text{append}((c, \text{old\_score} - \text{new\_score}))$
    \EndIf
\EndFor
\State Sort $\text{suggestions}$ by score (descending)
\State \Return $[\text{name for (name, score) in suggestions}]$
\end{algorithmic}
\end{algorithm}

\subsection{User Interface Workflow}
\label{subsec:cofactor-ui}

\begin{enumerate}
    \item User creates transition: \ce{Glucose -> G6P}
    \item System parses formulas, detects imbalance
    \item System suggests: ``Reaction imbalanced. Suggested cofactors: ATP (reactant), ADP (product). Apply?''
    \item User clicks ``Apply'' $\to$ System adds ATP input arc, ADP output arc
    \item User can fine-tune stoichiometry if needed
\end{enumerate}

\paragraph{Benefit:} Reduces manual entry errors, speeds up model construction.

%************************************************
\section{Applications and Benefits}
\label{sec:formula-applications}
%************************************************

\subsection{Model Validation}
\label{subsec:model-validation}

\paragraph{Use case:} User imports pathway from KEGG, converts to Extended Bio-PN.

\paragraph{Validation steps:}
\begin{enumerate}
    \item \textbf{Fetch formulas} from KEGG for all compounds
    \item \textbf{Verify elemental balance} for all reactions
    \item \textbf{Report imbalances}: ``Reaction R03270 (PFK) has \ce{H} imbalance: +1. Missing cofactor?''
    \item \textbf{Suggest fixes}: ``Add \ce{H2O} to reactants'' or ``Reaction proceeds via \ce{H+} release (buffered)''
\end{enumerate}

\paragraph{Result:} High-confidence model with verified stoichiometry.

\subsection{Pathway Completion}
\label{subsec:pathway-completion}

\paragraph{Use case:} User specifies partial pathway, wants to find missing reactions.

\paragraph{Algorithm:}
\begin{enumerate}
    \item Compute \textbf{elemental balance} for entire pathway
    \item Identify \textbf{accumulated elements}: $\sum S_e \cdot v \neq 0$ for some elements
    \item \textbf{Search reaction database} for reactions that consume/produce accumulated elements
    \item \textbf{Suggest candidate reactions}: ``Pathway accumulates 2 NADH. Consider adding NADH dehydrogenase.''
\end{enumerate}

\paragraph{Example (glycolysis):}
\begin{itemize}
    \item Reactions 1--10 convert \ce{Glucose -> 2 Pyruvate}
    \item Elemental balance: \ce{C}: 0, \ce{H}: --4, \ce{O}: 0, \ce{N}: 0, \ce{P}: 0
    \item \textbf{\ce{H} deficit}: 4 \ce{H} consumed (actually transferred to 2 NADH)
    \item Suggestion: ``Add \ce{NADH -> NAD+} reactions (e.g., lactate dehydrogenase, respiratory chain)''
\end{itemize}

\subsection{Redox Balance}
\label{subsec:redox-balance}

\paragraph{Extended feature:} Track \textbf{electron transfer} in redox reactions.

\paragraph{Example:} \ce{NAD+ + 2H+ + 2e- -> NADH + H+}

\paragraph{Elemental change:}
\begin{align}
\ce{NAD+} &: \ce{C21H27N7O14P2} \\
\ce{NADH} &: \ce{C21H29N7O14P2} \\
\Delta\ce{H} &= +2 \quad \text{(gained 2 \ce{H+})} \\
\text{Electrons} &: +2 e^- \quad \text{(reduced, gained electrons)}
\end{align}

\paragraph{Redox balance matrix} $S_{\text{redox}}$: Additional row tracking electrons.

\paragraph{Applications:}
\begin{itemize}
    \item Verify redox balance in respiratory chain
    \item Compute ATP/O ratio (ATP produced per oxygen reduced)
    \item Model mitochondrial electron transport
\end{itemize}

%************************************************
\section{Implementation Details}
\label{sec:formula-implementation}
%************************************************

\subsection{Data Structure}
\label{subsec:formula-data-structure}

\paragraph{Place attributes:}
\begin{itemize}
    \item \texttt{id}: Unique identifier
    \item \texttt{name}: Human-readable name
    \item \texttt{formula}: Elemental map (e.g., $\{\ce{C}: 6, \ce{H}: 12, \ce{O}: 6\}$)
    \item \texttt{formula\_string}: Hill notation (e.g., ``\ce{C6H12O6}'')
    \item \texttt{kegg\_id}: KEGG compound ID (e.g., ``C00031'')
    \item \texttt{chebi\_id}: ChEBI ID (e.g., ``CHEBI:17234'')
    \item \texttt{charge}: Net charge (e.g., 0, --2)
    \item \texttt{track\_protonation}: Boolean (True if \ce{H} balance critical)
\end{itemize}

\paragraph{Transition attributes:}
\begin{itemize}
    \item \texttt{id}: Unique identifier
    \item \texttt{name}: Human-readable name
    \item \texttt{reaction\_formula}: Reaction string (e.g., ``\ce{C6H12O6 + C10H16N5O13P3 -> C6H13O9P + C10H15N5O10P2 + H}'')
    \item \texttt{ec\_number}: Enzyme Commission number (e.g., ``2.7.1.1'' for Hexokinase)
    \item \texttt{reversible}: Boolean
    \item \texttt{kegg\_reaction\_id}: KEGG reaction ID (e.g., ``R00299'')
\end{itemize}

\subsection{Automatic Balance Checking}
\label{subsec:auto-balance-checking}

\paragraph{On model load/save:}

\begin{algorithm}[htbp]
\caption{Validate model elemental balance}
\label{alg:validate-model}
\begin{algorithmic}[1]
\Require Extended Bio-PN model
\Ensure List of warning messages (empty if all balanced)
\State $\text{warnings} \gets []$
\For{each transition $t$ in model}
    \State $\Delta \gets \texttt{verify\_elemental\_balance}(t)$
    \If{$\Delta \neq \emptyset$}
        \For{each element $e \in \Delta$}
            \If{$|\Delta[e]| > 2$ \textbf{and} $e \neq \ce{H}$} \Comment{Major imbalance}
                \State $\text{warnings}.\text{append}($``Transition $t$: \{$e$: $\Delta[e]$\}''$)$
            \ElsIf{$e = \ce{H}$ \textbf{and} $|\Delta[e]| > 5$} \Comment{Large \ce{H} imbalance}
                \State $\text{warnings}.\text{append}($``Transition $t$: \{\ce{H}: $\Delta[\ce{H}]$\} (check cofactors/water)''$)$
            \EndIf
        \EndFor
    \EndIf
\EndFor
\State \Return warnings
\end{algorithmic}
\end{algorithm}

%************************************************
\section{Limitations and Future Work}
\label{sec:formula-limitations}
%************************************************

\subsection{Current Limitations}
\label{subsec:current-limitations}

\subsubsection{Structural Isomers Not Distinguished}

\begin{itemize}
    \item G6P and F6P have same formula (\ce{C6H13O9P}) but different structures
    \item Elemental balance cannot detect ``fake'' isomerization errors
    \item \textbf{Solution}: Use SMILES or InChI for structural representation (future extension)
\end{itemize}

\subsubsection{Stereochemistry Not Captured}

\begin{itemize}
    \item D-glucose vs.\ L-glucose (mirror images, same formula)
    \item Enzymatic specificity requires structural information
    \item \textbf{Solution}: Annotate stereo centers, use SMILES with chirality
\end{itemize}

\subsubsection{Protonation State Ambiguity}

\begin{itemize}
    \item Different databases use different conventions (neutral, ionic, pH-dependent)
    \item Can cause $\pm$ few \ce{H} imbalances
    \item \textbf{Solution}: Standardize on KEGG neutral formulas, track charge separately
\end{itemize}

\subsection{Future Extensions}
\label{subsec:future-extensions}

\subsubsection{SMILES/InChI Integration}

\begin{itemize}
    \item Store structural formulas (SMILES) alongside elemental formulas
    \item Enable structural validation (e.g., bond rearrangement checks)
    \item Connect to cheminformatics tools (RDKit)
\end{itemize}

\subsubsection{Thermodynamic Integration}

\begin{itemize}
    \item Fetch $\Delta G^{\circ'}$ from eQuilibrator
    \item Compute $\Delta G$ based on current concentrations
    \item Enforce thermodynamic feasibility (disable reactions if $\Delta G \gg 0$)
\end{itemize}

\subsubsection{Isotope Tracking}

\begin{itemize}
    \item $^{13}$C-labeled glucose for metabolic flux analysis
    \item Track isotope distribution through pathways
    \item Enable isotope balance equations
\end{itemize}

%************************************************
\section{Summary}
\label{sec:chapter6-summary}
%************************************************

\textbf{Chapter 6} presented \textbf{biochemical formula tracking}, the fourth core innovation:

\paragraph{Key Contributions:}
\begin{enumerate}
    \item \textbf{Formalism}: Elemental composition map $\rho: P \to \text{BiochemicalFormula}$, reaction formulas $\rho: T \to \text{ReactionFormula}$
    \item \textbf{Balance verification}: Automatic checking that $\sum(\text{atoms}_{\text{in}}) = \sum(\text{atoms}_{\text{out}})$
    \item \textbf{Elemental balance matrix} $S_e$: Captures elemental stoichiometry, steady-state condition $S_e \cdot v = 0$
    \item \textbf{Database integration}: Fetch formulas from KEGG, ChEBI, PubChem
    \item \textbf{Cofactor suggestion}: Auto-detect imbalances, suggest missing cofactors
\end{enumerate}

\paragraph{Impact:}
\begin{itemize}
    \item Ensures biochemical validity (no phantom atoms)
    \item Reveals missing reactions (pathway completion)
    \item Speeds up model construction (auto-suggest cofactors)
    \item Enables advanced analyses (flux balance, redox balance)
\end{itemize}

\paragraph{Next Chapter:} Chapter 7 presents the \textbf{SBML import} system for automated model construction from systems biology repositories.
